% =============================================================================
% ISI Country Brief — Cyprus
% External Supplier Concentration Profile
%
% ISI Paper Series — Country Brief
% International Sovereignty Institute — Research Unit
% March 2026
% =============================================================================
\documentclass[12pt,a4paper]{article}

% --- ISI Paper Series shared template ----------------------------------------
\usepackage{isi-series}

% --- PDF metadata ------------------------------------------------------------
\hypersetup{
  pdftitle={External Supplier Concentration Profile: Cyprus — ISI EU-27 2024},
  pdfauthor={International Sovereignty Institute -- Research Unit},
  pdfcreator={International Sovereignty Institute},
  pdfsubject={ISI Paper Series -- Country Brief CY, EU-27, Vintage 2024},
  pdfkeywords={sovereignty index, EU-27, supplier concentration,
    cyprus, country brief},
}

% ==============================================================================
\begin{document}
\hypersetup{pageanchor=false}

% ---------- Title Page with Metadata ------------------------------------------
\begin{titlepage}
\centering
\vspace*{2cm}
{\Large\textsc{ISI Paper Series --- Country Brief}}\\[1.2cm]
{\LARGE\bfseries External Supplier Concentration Profile:\\[0.3em]
Cyprus}\\[1.5cm]
{\large International Sovereignty Index (ISI)\\[0.3em]
EU-27 --- Vintage 2024 --- Methodology v1.0}\\[1.2cm]
{\large International Sovereignty Institute --- Research Unit}\\[0.8cm]
{\large March 2026}

\vspace{0.8cm}
\noindent\rule{\textwidth}{0.4pt}
\vspace{0.4cm}

% --- Research Metadata --------------------------------------------------------
\begin{flushleft}
\noindent\textbf{Data basis:} ISI Paper~2 --- EU-27 Empirical Results 2024 (Methodology~v1.0, Vintage~2024).\\[4pt]
\noindent\textbf{Methodology:} ISI Paper~1 --- Methodological Foundations, Version~1.0.

\medskip

\noindent\textbf{JEL Codes:} F02, F15, F52, O30, Q43\\[4pt]
\noindent\textbf{Keywords:} sovereignty index; EU-27; supplier concentration; composite indicator; variance decomposition; defence dependence

\medskip

\noindent\textbf{Related publications in the ISI Paper Series:}
\begin{itemize}[nosep,leftmargin=1.5em]
\item ISI Paper~1 --- Technical Specification (v1.0) --- doi:\\
  \url{https://doi.org/10.5281/zenodo.18764227}
\item ISI Paper~2 --- EU-27 Empirical Results 2024 (v1.0) --- doi:\\
  \url{https://doi.org/10.5281/zenodo.18764170}
\end{itemize}
\end{flushleft}
\vfill
\end{titlepage}

\hypersetup{pageanchor=true}
\pagenumbering{arabic}
\pagestyle{fancy}
\fancyhead[L]{\small\sffamily\textit{ISI Country Brief}}
\fancyhead[R]{\small\sffamily\textit{Cyprus}}
\renewcommand{\headrulewidth}{0.4pt}

% ---------- Body --------------------------------------------------------------

\section{Executive Summary}
\label{sec:summary}

Cyprus ranks~2nd in the EU-27 with a composite \ISI score of~0.468, classified as \emph{moderately concentrated}.  The dominant axes are logistics~(0.941) and defence~(0.749).  The structural profile is characterised as logistics-led.  Under Dirichlet weight perturbation (10\,000~Monte Carlo draws), Cyprus' mean rank is~2.11 with a standard deviation of~0.33; the 5th--95th percentile band spans ranks~2 to~3.


\section{Country Position in the EU-27}
\label{sec:position}

Cyprus occupies rank~2 of~27 EU member states (composite score~0.468).  The EU-27 mean composite score is~0.344 with a standard deviation of~0.070.  Cyprus' score lies above the EU-27 mean, indicating above-average external supplier concentration.  In the ranking, Cyprus is situated between Malta (rank~1; 0.518) and Ireland (rank~3; 0.428).


\section{Axis Profile}
\label{sec:axis-profile}

\Cref{tab:axis-profile} presents Cyprus' scores on all six \ISI axes.

\begin{table}[htbp]
\centering\small
\caption{Cyprus' \ISI axis profile.}
\label{tab:axis-profile}
\begin{tabular}{@{}lr@{}}
\toprule
\textbf{Axis} & \textbf{Score} \\
\midrule
    Finance           & 0.124 \\
    Energy            & 0.348 \\
    Technology        & 0.224 \\
    Defence           & 0.749 \\
    Critical Inputs   & 0.424 \\
    Logistics         & 0.941 \\
\bottomrule
\end{tabular}
\end{table}

The highest axis score is logistics~(0.941), which lies above the EU-27 axis mean of~0.518.  The second-highest axis is defence~(0.749), above the EU-27 axis mean of~0.573.  The lowest axis score is finance~(0.124), below the EU-27 mean of~0.148.


\section{Structural Drivers}
\label{sec:drivers}

\begin{sloppypar}
Cyprus' composite score is led by logistics~(0.941), with defence~(0.749) as the second contributor.  The gap between the two highest axes is~0.192.  While logistics exerts the strongest influence on the composite, the contribution of defence is non-negligible.
\end{sloppypar}


\section{Comparative Context}
\label{sec:comparative}

Across the EU-27, the covariance-based variance decomposition shows that defence (35.1\,\%) and logistics (34.3\,\%) together account for 69.4\,\% of total cross-sectional composite variance.  Cyprus' defence score~(0.749) lies above the EU-27 defence mean~(0.573), and its logistics score~(0.941) lies above the EU-27 logistics mean~(0.518).  Under Dirichlet weight perturbation, Cyprus' rank standard deviation is~0.33, indicating high rank stability.



Cyprus occupies a distinct position in the upper segment of the EU-27 distribution.  The composite gap to Malta~(rank~1) is~0.050, substantially larger than the gap to Ireland~(rank~3; 0.040).  This asymmetry suggests that while Cyprus is clearly separated from the top-ranked outlier, its distance from the countries immediately below is comparatively modest.

The axis profile exhibits a pronounced gradient: the two highest axes---defence~(0.749) and logistics~(0.941)---are both well above their EU-27 means, while the four remaining axes cluster in the range 0.124--0.424.  This bimodal structure aligns with the EU-wide variance decomposition, in which defence and logistics jointly account for 69.4\,\% of composite variance.  Cyprus's rank position is therefore shaped predominantly by the axes along which cross-country differentiation is greatest.

The 5th--95th~percentile rank band of~2 to~3 under Dirichlet perturbation is narrow, reflecting the sizeable composite distance from countries in the mid-ranking.  Nevertheless, the band does extend to rank~3, indicating that under certain weight configurations the ordinal boundary between Cyprus and Ireland becomes marginal.


\section{Interpretation Boundary}
\label{sec:interpretation}

The \ISI measures concentration of external suppliers, not sovereignty in a normative or political sense.  High concentration may reflect geography, economic specialisation, or alliance structures.  A high composite score does not imply policy failure; a low score does not imply policy success.  The index provides an empirical measurement basis --- what policymakers derive from it is a separate question.


% ---------- References --------------------------------------------------------
\clearpage
\vspace*{1.5cm}

\begingroup\emergencystretch=2em\hbadness=10000
\section*{References}

\begin{enumerate}[label={[\arabic*]},leftmargin=2em,itemsep=6pt]
\item International Sovereignty Institute.
  \textit{International Sovereignty Index (ISI) technical specification, version~1.0}.
  Zenodo, 2026.\\ISI Paper Series, Paper~1.
  doi:\,\href{https://doi.org/10.5281/zenodo.18764227}{10.5281/zenodo.18764227}.

\item International Sovereignty Institute.
  \textit{International Sovereignty Index (ISI): EU-27 empirical results 2024 (v1.0)}.
  Zenodo, 2026.\\ISI Paper Series, Paper~2.
  doi:\,\href{https://doi.org/10.5281/zenodo.18764170}{10.5281/zenodo.18764170}.
\end{enumerate}
\endgroup

\end{document}
